\section{Case studies и пространственная демонстрация}

\subsection{Роль case studies в работе}

После построения предиктивного ядра, перехода к рекомендательной постановке, confidence-анализа и explainability-слоя возникает задача показать, как система ведёт себя не только на уровне усреднённых метрик, но и на уровне отдельных наблюдений. Для этого в работе используются case studies --- специально отобранные примеры полей, на которых можно последовательно проследить связь между историей культур, результатом модели, уровнем уверенности и графовым интерпретационным контекстом.

Такой формат анализа важен по двум причинам. Во-первых, усреднённые метрики не показывают, как именно shortlist работает в простых и сложных случаях. Во-вторых, explainability-слой по своей природе является локальным: его смысл раскрывается не на уровне общей таблицы результатов, а на уровне конкретного списка кандидатов для конкретного поля. Поэтому case studies в данной работе выступают как естественный способ показать, каким образом предиктивный, confidence- и graph-based слои взаимодействуют между собой.

В проекте были собраны четыре case studies:

\begin{itemize}
\tightlist
\item
  \texttt{high\ confidence};
\item
  \texttt{medium\ confidence};
\item
  \texttt{low\ confidence};
\item
  \texttt{difficult\ /\ disagreement\ case}.
\end{itemize}

Этот набор позволяет охватить как относительно простые и устойчивые сценарии, так и случаи, где shortlist требует более внимательной интерпретации.

\subsection{High confidence case}

Первый тип case studies соответствует ситуациям, в которых модель демонстрирует высокий уровень уверенности, а top-1 кандидат выглядит содержательно согласованным с explainability-контекстом. В таких случаях история поля, вероятностный вывод модели и graph-based support обычно не противоречат друг другу. Это наиболее простой сценарий работы системы: shortlist остаётся полезным, но первая позиция уже выглядит достаточно устойчивой.

С точки зрения интерпретации такой кейс важен потому, что показывает нормальное поведение системы в условиях относительно «чистой» структуры данных. Если поле относится к хорошо распознаваемому паттерну, модель формирует сильный top-1 прогноз, а explainability-слой подтверждает, что выбранный кандидат согласуется с переходной статистикой, региональным контекстом и активными правилами. Для пользователя это означает, что система не просто выдаёт высокий score, но и делает это в достаточно понятной и интерпретируемой ситуации.

\subsection{Medium confidence case}

Второй тип case studies соответствует случаям средней уверенности. Здесь top-1 кандидат остаётся содержательно правдоподобным, однако shortlist в целом становится более важным, поскольку помимо первой позиции появляются и другие кандидаты, которые могут выглядеть осмысленно с точки зрения истории поля и explainability-контекста.

Такой кейс особенно важен для прикладной логики работы. Он показывает, что рекомендательная постановка полезна не только для обработки явных и простых ситуаций, но и для анализа наблюдений, где система не должна ограничиваться одним ответом. В подобных случаях shortlist играет роль пространства разумных альтернатив: модель указывает основное направление, но explainability-слой позволяет рассматривать и соседние варианты как содержательно поддержанные кандидаты.

\subsection{Low confidence case}

Третий тип case studies связан со случаями низкой уверенности. Именно такие наблюдения особенно хорошо показывают, почему confidence-слой и Top-k постановка являются принципиально важными для всей системы. В low-confidence сценариях top-1 рекомендация может быть недостаточно устойчивой, однако shortlist в целом остаётся полезным, поскольку в нём по-прежнему присутствуют правдоподобные альтернативы.

На уровне explainability это проявляется в том, что graph-based context может поддерживать не только top-1 кандидата, но и одну или несколько альтернатив. Для пользователя это означает, что результат не следует воспринимать как сильное однозначное предписание. Напротив, низкая уверенность должна побуждать рассматривать несколько кандидатов одновременно и использовать дополнительный предметный анализ. Такой кейс хорошо демонстрирует, что при слабом модельном сигнале shortlist-подход остаётся более устойчивым форматом вывода, чем жёсткий одновариантный ответ.

\subsection{Difficult / disagreement case}

Четвёртый тип case studies соответствует наиболее сложным ситуациям, в которых наблюдается напряжение между top-1 выбором модели и explainability-контекстом. Именно такие кейсы позволяют лучше всего показать, зачем в архитектуре вообще нужен graph-based слой. В этих сценариях первая позиция shortlist может не совпадать с наиболее сильно поддержанным графом кандидатом, а фактическая следующая культура может находиться не на первом месте, а среди альтернатив внутри Top-k списка.

Интерпретация таких случаев является особенно важной. Они не означают, что модель бесполезна или что граф «лучше» модели. Их смысл заключается в другом: shortlist следует рассматривать как пространство конкурирующих вариантов, где первая позиция не всегда исчерпывает всю полезную информацию. Если explainability-слой сильнее поддерживает альтернативного кандидата, это указывает на спорность ситуации и делает case-level анализ более содержательным. Для задачи поддержки принятия решений такой результат даже важнее, чем ещё один простой случай правильного top-1 прогноза, поскольку он показывает, как система ведёт себя в неоднозначных наблюдениях.

\subsection{Общая интерпретация case studies}

Рассмотренные case studies позволяют сделать несколько важных выводов. Во-первых, они подтверждают, что система должна интерпретироваться не как одновариантный классификатор, а как рекомендательный прототип, полезность которого особенно проявляется в анализе нескольких кандидатов одновременно. Во-вторых, они показывают, что confidence-слой действительно задаёт различие между более устойчивыми и более спорными рекомендациями. В-третьих, они демонстрируют, что knowledge graph полезен не как замена модели, а как механизм интерпретации shortlist и сравнения кандидатов внутри него.

Таким образом, case studies в работе выполняют не иллюстративную, а аналитическую функцию. Они связывают между собой все уровни системы:

\begin{itemize}
\tightlist
\item
  историю поля;
\item
  модельный shortlist;
\item
  confidence-сигнал;
\item
  graph-based explanation.
\end{itemize}

\subsection{Пространственная демонстрация}

Следующим шагом после case-level анализа становится пространственная демонстрация результатов. Её задача состоит в том, чтобы связать рекомендации с конкретными полями и показать, как итоговый результат может быть представлен в пользовательском интерфейсе. Пространственный слой не создаёт новой модели и не влияет на метрики качества. Он выполняет демонстрационную и аналитическую функцию: переносит recommendation и explainability-результаты в карту полей.

В проекте были сформированы несколько spatial artifacts, в том числе:

\begin{itemize}
\tightlist
\item
  \texttt{spatial\_explainable\_demo\_map.html};
\item
  \texttt{spatial\_bbox\_overview\_map.html};
\item
  \texttt{spatial\_bbox\_demo\_map.html};
\item
  \texttt{spatial\_layer\_comparison\_summary.csv};
\item
  \texttt{spatial\_bbox\_selection\_summary.csv}.
\end{itemize}

Эти артефакты используются для того, чтобы показать, как пользователь может перейти от общего spatial context к анализу конкретного поля. В итоговом варианте карта связывает три типа информации:

\begin{itemize}
\tightlist
\item
  геометрическое или пространственное представление поля;
\item
  recommendation output для этого поля;
\item
  explainability-контекст для выбранного кейса.
\end{itemize}

Иными словами, spatial demo позволяет перевести результат из формата таблиц и локальных графов в более прикладной сценарий визуального анализа.

\subsection{Полный spatial context и demo subset}

На одном из этапов пространственного анализа было важно понять, не выглядит ли демонстрационная карта искусственно разреженной из-за того, что для показа отбирается только небольшой набор case studies. Для этого был выполнен отдельный анализ spatial layers. Он показал, что полный региональный слой содержит 19 919 полей, а после фильтрации до рабочего пространства из 9 целевых классов сохраняется 16 663 поля, то есть около 83.7\% полного слоя. При этом сам демонстрационный subset включает лишь 20 полей.

Этот результат важен для интерпретации spatial demo. Он показывает, что визуальная «разреженность» демо-кейсов связана не с тем, что фильтрация классов разрушила spatial структуру региона, а с тем, что на карте выделяется очень маленький curated subset для explainability-анализа. В финальной логике демонстрации это привело к двухслойной подаче:

\begin{itemize}
\tightlist
\item
  фоновый слой показывает полный filtered spatial context;
\item
  поверх него выделяются demo cases, используемые для подробного анализа.
\end{itemize}

Такое решение делает карту более содержательной: пользователь видит, что рассматриваемые поля принадлежат плотному сельскохозяйственному региону, а выбранные кейсы являются не случайными точками, а подмножеством общего spatial context.

\subsection{Bounding box selection}

Для повышения читаемости и удобства демонстрации в работу был добавлен отдельный этап локального spatial selection через bounding box. Его задача состоит в том, чтобы ограничить карту компактной областью, внутри которой можно удобно анализировать поля и связанные с ними рекомендации. Такой подход полезен не только для презентации, но и для самой логики spatial demo: пользователь работает не со всем регионом сразу, а с локальным фрагментом, где поля визуально сгруппированы и доступны для детального просмотра.

По данным \texttt{spatial\_bbox\_selection\_summary.csv}, внутри выбранного bbox-контекста содержится 157 полей, из которых с test-строками сопоставлено 23 поля; итоговый clickable demo dataset также содержит 23 поля. Это означает, что bounding box selection позволяет перейти от общего spatial context к более компактному и управляемому пространственному подмножеству, не теряя при этом связки между полем, рекомендацией и explainability-контекстом.

С архитектурной точки зрения bounding box selection не меняет саму систему рекомендаций. Он лишь добавляет spatial filtering mechanism, который делает демонстрацию более наглядной и ближе к реальному пользовательскому сценарию. Благодаря этому пространственная часть работы воспринимается не как отдельный GIS-эксперимент, а как естественное продолжение explainable recommendation prototype.

\subsection{Место spatial demo в общей работе}

Пространственная демонстрация завершает прикладную часть исследования. Если предыдущие главы были сосредоточены на данных, модели, recommendation-логике, confidence и explainability, то spatial demo показывает, как эти результаты могут быть представлены в интерфейсе, ориентированном на работу с конкретными полями. Это особенно важно для дипломной работы, поскольку позволяет продемонстрировать не только качество модели как таковой, но и форму её практического использования.

При этом spatial demo не изменяет основной научный вывод работы. Финальным предиктивным ядром остаётся baseline CatBoost, а карта и bbox-selection выступают как интерфейсно-аналитические надстройки. Их роль заключается в том, чтобы сделать recommendation и explainability-результаты более наглядными и пригодными для визуального анализа.

\subsection{Выводы по главе}

В данной главе были рассмотрены case studies и пространственная демонстрация результатов. Case studies показали, как система ведёт себя в различных типах наблюдений --- от простых high-confidence случаев до сложных disagreement-сценариев. Это подтвердило, что полезность системы заключается не только в top-1 предсказании, но и в возможности анализировать shortlist кандидатов с учётом confidence- и graph-based сигналов.

Пространственная демонстрация дополнила эту логику визуальным уровнем представления. Она позволила связать рекомендации с геометрией полей, выделить demo cases на фоне полного filtered spatial context и организовать локальный анализ через bounding box selection. В результате итоговая система может рассматриваться не только как модель и набор метрик, но и как целостный explainable recommendation prototype, для которого предусмотрен наглядный сценарий пользовательской работы.
